% GENERATED FILE. Edit canonical lessons, then run npm run generate:books.
% canonical-source-hash: fnv1a64:4f440bb1e751654b
% canonical-lessons: FR-C508-coureur, FR-C508-ligne-d-arrivee, FR-C508-marathon, FR-C508-relais, FR-C508-chronometrer

\chapter{La course}
\label{ch:fr-la-course}
\hlchaptermodality{508}

\begin{chapteropening}
\textbf{By the end of this chapter:} \emph{I can talk about races and runners.}

Complete the last lesson of chapter 508: I can talk about races and runners.
\end{chapteropening}

\section[le coureur]{le coureur — runner; rider (in a race)}

\label{lesson:FR-C508-coureur}



\begin{quote}
Before the new one: say the French for champion, then the French for to face.
\end{quote}

\subsection*{You'll want to know: le coureur}
\textbf{le coureur} — \textquotedblleft{}runner; rider (in a race)\textquotedblright{}.

\begin{quote}
Le coureur a battu son propre record.
\end{quote}

\begin{quote}
\emph{The runner beat his own record.}
\end{quote}

A woman is \emph{la coureuse}. Also a cyclist in a race: \emph{un coureur cycliste}.

\begin{cousinweb}
\emph{le coureur} is built on \emph{courir} \textquotedblleft{}to run\textquotedblright{}, inherited from Latin \emph{currere} \textquotedblleft{}to run\textquotedblright{} (Lewis \& Short).
\end{cousinweb}

\begin{grammarlens}[title={What you've built}]
One more word for telling, explaining and describing at B1.
\end{grammarlens}

\subsection*{Your turn}
\begin{itemize}
\raggedright
  \item \emph{Say it:} \emph{le coureur}
  \item \emph{Say it:} \emph{le coureur}, once more
  \item \emph{Say it:} a sentence of your own with \emph{le coureur}
  \item \emph{Recall:} say the French for champion, then the French for to face, then say \emph{le coureur} again
\end{itemize}

\begin{tcolorbox}[breakable,colback=teal!4,colframe=teal!35!black,title={Before you move on}]
What does le coureur mean? (\textquotedblleft{}Runner; rider (in a race)\textquotedblright{}.) Say it once more.
\end{tcolorbox}
\section[la ligne d'arrivée]{la ligne d'arrivée — finishing line}

\label{lesson:FR-C508-ligne-d-arrivee}



\begin{quote}
Before the new one: say the French for to face, then the French for runner.
\end{quote}

\subsection*{You'll want to know: la ligne d'arrivée}
\textbf{la ligne d'arrivée} — \textquotedblleft{}finishing line\textquotedblright{}.

\begin{quote}
Elle a franchi la ligne d'arrivée en premier.
\end{quote}

\begin{quote}
\emph{She crossed the finishing line first.}
\end{quote}

The start is \emph{la ligne de départ}.

\begin{cousinweb}
\emph{la ligne d'arrivée} joins \emph{ligne}, from Latin \emph{līnea} \textquotedblleft{}a line\textquotedblright{}, and \emph{arrivée}, built on \emph{arriver}, from Vulgar Latin \emph{arrīpāre} \textquotedblleft{}to come to the shore\textquotedblright{}, built on \emph{rīpa} \textquotedblleft{}bank\textquotedblright{} (Lewis \& Short).
\end{cousinweb}

\begin{grammarlens}[title={What you've built}]
One more word for telling, explaining and describing at B1.
\end{grammarlens}

\subsection*{Your turn}
\begin{itemize}
\raggedright
  \item \emph{Say it:} \emph{la ligne d'arrivée}
  \item \emph{Say it:} \emph{la ligne d'arrivée}, once more
  \item \emph{Say it:} a sentence of your own with \emph{la ligne d'arrivée}
  \item \emph{Recall:} say the French for to face, then the French for runner, then say \emph{la ligne d'arrivée} again
\end{itemize}

\begin{tcolorbox}[breakable,colback=teal!4,colframe=teal!35!black,title={Before you move on}]
What does la ligne d'arrivée mean? (\textquotedblleft{}Finishing line\textquotedblright{}.) Say it once more.
\end{tcolorbox}
\section[le marathon]{le marathon — marathon}

\label{lesson:FR-C508-marathon}



\begin{quote}
Before the new one: say the French for runner, then the French for finishing line.
\end{quote}

\subsection*{You'll want to know: le marathon}
\textbf{le marathon} — \textquotedblleft{}marathon\textquotedblright{}.

\begin{quote}
Mon père a couru son premier marathon à cinquante ans.
\end{quote}

\begin{quote}
\emph{My father ran his first marathon at fifty.}
\end{quote}

Also a long, tiring task: \emph{un marathon de réunions}.

\begin{cousinweb}
\emph{le marathon} takes its name from a Greek town (EtymDB), whose own name comes from Greek \emph{márathon} \textquotedblleft{}fennel\textquotedblright{}; the race recalls a runner's long journey from there.
\end{cousinweb}

\begin{grammarlens}[title={What you've built}]
One more word for telling, explaining and describing at B1.
\end{grammarlens}

\subsection*{Your turn}
\begin{itemize}
\raggedright
  \item \emph{Say it:} \emph{le marathon}
  \item \emph{Say it:} \emph{le marathon}, once more
  \item \emph{Say it:} a sentence of your own with \emph{le marathon}
  \item \emph{Recall:} say the French for runner, then the French for finishing line, then say \emph{le marathon} again
\end{itemize}

\begin{tcolorbox}[breakable,colback=teal!4,colframe=teal!35!black,title={Before you move on}]
What does le marathon mean? (\textquotedblleft{}Marathon\textquotedblright{}.) Say it once more.
\end{tcolorbox}
\section[le relais]{le relais — relay; relay race}

\label{lesson:FR-C508-relais}



\begin{quote}
Before the new one: say the French for finishing line, then the French for marathon.
\end{quote}

\subsection*{You'll want to know: le relais}
\textbf{le relais} — \textquotedblleft{}relay; relay race\textquotedblright{}.

\begin{quote}
Notre équipe a gagné le relais quatre fois cent mètres.
\end{quote}

\begin{quote}
\emph{Our team won the four by one hundred metres relay.}
\end{quote}

\emph{Prendre le relais} — to take over.

\begin{cousinweb}
\emph{le relais} comes from Old French \emph{relaier} \textquotedblleft{}to leave behind\textquotedblright{} (EtymDB), built from \emph{re-} and \emph{laier} \textquotedblleft{}to leave\textquotedblright{}: fresh horses were left waiting along the road.
\end{cousinweb}

\begin{grammarlens}[title={What you've built}]
One more word for telling, explaining and describing at B1.
\end{grammarlens}

\subsection*{Your turn}
\begin{itemize}
\raggedright
  \item \emph{Say it:} \emph{le relais}
  \item \emph{Say it:} \emph{le relais}, once more
  \item \emph{Say it:} a sentence of your own with \emph{le relais}
  \item \emph{Recall:} say the French for finishing line, then the French for marathon, then say \emph{le relais} again
\end{itemize}

\begin{tcolorbox}[breakable,colback=teal!4,colframe=teal!35!black,title={Before you move on}]
What does le relais mean? (\textquotedblleft{}Relay; relay race\textquotedblright{}.) Say it once more.
\end{tcolorbox}
\section[chronométrer]{chronométrer — to time (with a stopwatch)}

\label{lesson:FR-C508-chronometrer}



\begin{quote}
Before the new one: say the French for marathon, then the French for relay.
\end{quote}

\subsection*{You'll want to know: chronométrer}
\textbf{chronométrer} — \textquotedblleft{}to time (with a stopwatch)\textquotedblright{}.

\begin{quote}
L'entraîneur chronomètre chaque coureur.
\end{quote}

\begin{quote}
\emph{The coach times each runner.}
\end{quote}

\emph{Un chronomètre}, informally \emph{un chrono} — a stopwatch.

\begin{cousinweb}
\emph{chronométrer} is built on \emph{chronomètre} \textquotedblleft{}stopwatch\textquotedblright{}, a learned word made from Greek \emph{khrónos} \textquotedblleft{}time\textquotedblright{} and \emph{métron} \textquotedblleft{}measure\textquotedblright{}.
\end{cousinweb}

\begin{grammarlens}[title={What you've built}]
One more word for telling, explaining and describing at B1.
\end{grammarlens}

\subsection*{Your turn}
\begin{itemize}
\raggedright
  \item \emph{Say it:} \emph{chronométrer}
  \item \emph{Say it:} \emph{chronométrer}, once more
  \item \emph{Say it:} a sentence of your own with \emph{chronométrer}
  \item \emph{Recall:} say the French for marathon, then the French for relay, then say \emph{chronométrer} again
\end{itemize}

\begin{tcolorbox}[breakable,colback=teal!4,colframe=teal!35!black,title={Before you move on}]
What does chronométrer mean? (\textquotedblleft{}To time (with a stopwatch)\textquotedblright{}.) Say it once more.
\end{tcolorbox}
